We study how three action heads---an MSE-regression MLP, a Gaussian-mixture head and a 50-step DDPM head, all reimplemented here---cope with bimodal demonstrations in a 2D point-mass reach task in which an obstacle can be passed on either side. With 2,000 scripted demonstrations, 200 closed-loop rollouts and 10 seeds, the MSE head reaches the goal in only \rhval{main/mse-mlp/bimodal50/success/mean:pct1}\% of rollouts (the rest collide with the obstacle) when start states are nearly identical, whereas the mixture and the DDPM head reach \rhval{main/gmm-head/bimodal50/success/mean:pct1}\% and \rhval{main/ddpm-head/bimodal50/success/mean:pct1}\%. The DDPM head is not distinguishable from a two-component mixture in success (Welch $p=\rhval{cmp/main/gmm-head/bimodal50/success/welch_p:2}$) or mode share. The MSE failure is fragile: it depends on start-state jitter (success \rhval{sweep_jitter/mse-mlp-jit-0.0/bimodal50/success/mean:pct1}\% with exactly identical starts, \rhval{sweep_jitter/mse-mlp-jit-0.1/bimodal50/success/mean:pct1}\% with large jitter) and, on skewed demonstrations, MSE succeeds by collapsing to the majority mode. Two further ablations show where the diffusion head is weaker: a single denoising step fails and long open-loop action chunks hurt DDPM but not the mixture. This is a small CPU study in one toy environment; it supports a narrow claim about when multimodal heads are needed, not a ranking of generative policy classes.
